\documentclass{article}
\usepackage{amsmath,hyperref}
\begin{document}
\section*{A cross-supplier return bound to the same persisted receipt}
\paragraph{Compared interfaces.}
The native warehouse worker commits the old operation, then exits.
A separate export worker also exits, after persisting the closed
archive. The donor is moved out of the online working set. P and Q
prepare copies by SQLite backup to separate files, each reopened
read-only for the send. Their agreement is on the same named
old operation, source epoch, archive digest, and fixed codec. The
dispatcher reads a local grant table and requires the copy to match
both its grant and the seal actually received earlier at Base.
This recovery service does not have access to the donor or original
shared witness at the action cut. Separate files are not independent
failure domains, and two snapshots of the same witness are not
independent verification of its truth.

P's native header arrives before the relay change and its native early
result digit arrives at the next cut. The early digit is evidence
about this immutable record and codec position, not a delegation to
Q. A Q continuation is considered only if its own current grant and
persisted copy independently qualify, the early digit matches its
copy, and its route and transmit and receive budgets permit the
action. Generation and codec agreement precede byte joining. The
separate application still reconstructs the full archive, applies the
original digest and SQLite verification, activates the authorized
station and retries the named operation. A cache reply carrying the exact committed
warehouse tick and outcome establishes the measured settlement. Merely dispatching
does not.

\paragraph{An effect after the native pose event.}
Both arms get the same prepared copies and the same exact early
observations. The queued arm commits the eligible P and Q requests
before the new pose observation. The local admission arm defers
sender choice to the proposed serial seam after the native visibility
rebuild and before queued effects. It applies the existing admission
core to earned copy and route facts and the live counters, preferring
the old return with independent auxiliary work over the old return
alone. A direct same-cut rule agrees. The ordered seam is an
instrumented simulator prototype, not a stock API or a proof of
snapshot consistency under arbitrary concurrency.

With a late near registration, P's route is no longer usable and Q's
is. At prior debits $d=41,\ q=151$, the 164-byte continuation is
exactly affordable for Q after its two earlier charges. The old
P callback alone cannot establish what to do. The exact cross-supplier
return can use the early P digit followed by the Q remainder.
The blind queued comparison attempts P even though the route is no
longer usable. Its charged lost attempt consumes the bytes needed to
send the independent four-byte auxiliary report. The local policy
charges only Q and keeps the auxiliary work. With both routes
available it can still use Q to protect the auxiliary obligation.
An empty or revoked Q grant, a stale generation, or a changed receipt
key with the same old digit cannot be spliced. At a receive debit of
157, neither route can deliver the 164-byte remainder under the
320-byte shared cap; freshness of either route does not remove that
constraint.

\paragraph{Boundary.}
The two-copy and authorization setup, exact descriptor reuse, codec and
slot identification, and local reads are supplied by the declared
facility in both comparison arms, not claimed free to deploy or
established by the radio bit. There is no claim that copying is
live-fault-atomic, that late authorization revocation after the serial
cut is covered, or that the modeled RF law measures a physical
channel. The measured improvement is against fixed queued replication
at the declared interface. An ordinary scheduler with the same
qualified cut and access to dispatch ties. This is held-in evidence for
a compositional capability, not independent mission acceptance.
\end{document}
